\subsection{Topology of the Galois group}

Let N be a Galois extension of K and \(\Gamma\) the Galois group of N over K. We equip N with the discrete topology, the set \(N^{N}\) of all mappings of N into itself with the product topology of the discrete topology of the factors (« topology of simple convergence in N ») and the group \(\Gamma\) with the topology induced from \(N^{N}\) .

Let \(\Lambda\) be the set of all subextensions of \(\mathbf{N}\) of finite degree over \(\mathbf{K}\) . For \(\sigma \in \Gamma\) and \(\mathrm{E} \in \Lambda\) we shall write \(\mathrm{U}_{\mathrm{E}}(\sigma)\) for the set of elements \(\tau\) of \(\Gamma\) which have the same restriction as \(\sigma\) to \(\mathrm{E}\) . If \(\mathrm{E} = \mathrm{K}(x_1, \ldots, x_n)\) , the set \(\mathrm{U}_{\mathrm{E}}(\sigma)\) consists of the elements \(\tau \in \Gamma\) such that \(\tau(x_1) = \sigma(x_1), \ldots, \tau(x_n) = \sigma(x_n)\) . It follows that the family \((\mathrm{U}_{\mathrm{E}}(\sigma))_{\mathrm{E} \in \Lambda}\) is a base of the filter of neighbourhoods of \(\sigma\) in \(\Gamma\) .

When N is of finite degree over K, we have \(N \in \Lambda\) and \(U_{N}(\sigma) = \{\sigma\}\) , hence the topology of \(\operatorname{Gal}(N/K)\) is discrete; we recall (V, p.~58), that the group \(\operatorname{Gal}(N/K)\) is finite in this case.

This description of the topology of \(\operatorname{Gal}(\mathbf{N}/\mathbf{K})\) shows that the restriction homomorphism of \(\operatorname{Gal}(\mathbf{N}/\mathbf{K})\) onto \(\operatorname{Gal}(\mathbf{L}/\mathbf{K})\) is continuous for every subextension L of N which is Galois over N.

Let A be a subset of \(\Gamma\) . To say that A is open means that for each \(\sigma\in A\) there exists E in \(\Lambda\) such that the set \(U_{E}(\sigma)\) is contained in A. The closure \(\bar{A}\) of A consists of all \(\sigma\in\Gamma\) such that for any \(E\in\Lambda\) there exists \(\tau\in A\) having the same restriction to E as \(\sigma\) ; the field of invariants of \(\bar{A}\) is the same as that of A.

Let \(\varepsilon\) be the neutral element of \(\Gamma\) and let \(\Lambda'\) be the set of subextensions of \(N\) which are Galois and of finite degree over \(K\) . By Prop. 2 (V, p.~57) the set \(\Lambda'\) is cofinal in \(\Lambda\) and the family \((U_E(\varepsilon))_{E \in \Lambda'}\) is thus a base of the filter of neighbourhoods of \(\varepsilon\) in \(\Gamma\) . Further, for \(E \in \Lambda'\) the set \(U_E(\varepsilon)\) is the kernel of the restriction homomorphism of \(\operatorname{Gal}(N/K) = \Gamma\) into \(\operatorname{Gal}(E/K)\) . Since \(\operatorname{Gal}(E/K)\) is finite, it follows that \(U_E(\varepsilon)\) is an open and closed subgroup, normal and of finite index in \(\Gamma\) .

Clearly we have \(\mathrm{U}_{\mathrm{E}}(\sigma)=\sigma\mathrm{U}_{\mathrm{E}}(\varepsilon)=\mathrm{U}_{\mathrm{E}}(\varepsilon)\sigma\) for \(\sigma\in\Gamma\) and \(E\in\Lambda'\) . Since \(\mathrm{U}_{\mathrm{E}}(\varepsilon)\) is a normal subgroup of \(\Gamma\) for all \(E\in\Lambda'\) and the family \((\mathrm{U}_{\mathrm{E}}(\varepsilon))_{E\in\Lambda'}\) is a base of neighbourhoods at \(\varepsilon\) , the topology of \(\Gamma\) is compatible with the group structure (Gen.~Top., III, p.~223). In other words, the mapping \((\sigma,\tau)\mapsto\sigma\tau^{-1}\) of \(\Gamma\times\Gamma\) into \(\Gamma\) is continuous.

PROPOSITION 4. --- Let N be a Galois extension of K. Then the Galois group \(\Gamma = \operatorname{Gal}(\mathbf{N}/\mathbf{K})\) is compact and totally disconnected.

Every element \(\sigma\) of \(\Gamma\) has a base of neighbourhoods consisting of the open and closed sets \(U_{E}(\sigma)\) , hence \(\Gamma\) is totally disconnected (Gen.~Top., I, p.~111). We have \(\{\sigma\} = \bigcap_{E \in \Lambda} U_{E}(\sigma)\) , hence \(\Gamma\) is separated. For each \(x \in \mathbb{N}\) the set of conjugates

\(\sigma(x)\) of \(x\) , as \(\sigma\) ranges over \(\Gamma\) , is finite because \(x\) is algebraic over K (V, p.~53, Cor. 1); hence all the projections of \(\Gamma\) on the factor spaces of \(\mathbf{N}^{\mathbb{N}}\) are finite sets, and this shows \(\Gamma\) to be relatively compact in \(\mathbf{N}^{\mathbb{N}}\) (Gen.~Top., I, p.~88). It remains to show that \(\Gamma\) is closed in \(\mathbf{N}^{\mathbb{N}}\) . Now if \(u\) is in the closure of \(\Gamma\) in \(\mathbf{N}^{\mathbb{N}}\) , then for each pair of points \((x,y)\) of N there exists \(\sigma \in \Gamma\) with \(u(x) = \sigma(x)\) , \(u(y) = \sigma(y)\) , \(u(x + y) = \sigma(x + y)\) , \(u(xy) = \sigma(xy)\) , whence \(u(x + y) = u(x) + u(y)\) and \(u(xy) = u(x)u(y)\) . By the same reasoning we have \(u(x) = x\) for all \(x \in \mathbf{K}\) , hence \(u\) is a K-homomorphism of N into N; since N is algebraic over K, \(u\) is a K-automorphism of N (V, p.~52, Prop. 1), hence \(u \in \Gamma\) .

Let N be a Galois extension of K and \((\mathbf{N}_{i})_{i\in I}\) an increasing directed family of subextensions of N. Suppose that \(N_{i}\) is Galois over K for all \(i\in I\) and that \(N=\bigcup_{i\in I}N_{i}\) . For each \(i\in I\) denote by \(\Gamma_{i}\) the Galois group of \(N_{i}\) over K; for

\(i \leqslant j\) in I we have \(N_i \subset N_j\) and the restriction homomorphism \(\varphi_{ij}\) of \(\Gamma_j\) into \(\Gamma_i\) is defined. It is continuous and so the family \((\Gamma_i, \varphi_{ij})\) is an inverse system of topological groups. Further, for each \(i \in I\) denote by \(\lambda_i\) the restriction homomorphism of \(\text{Gal}(N/K)\) into \(\text{Gal}(N_i/K) = \Gamma_i\) ; it is continuous and we have \(\lambda_i = \varphi_{ij} \circ \lambda_j\) for \(i \leqslant j\) , hence the family \((\lambda_i)_{i \in I}\) defines a continuous homomorphism \(\lambda\) of \(\text{Gal}(N/K)\) into \(\varprojlim \Gamma_i\) .

PROPOSITION 5. --- The homomorphism \(\lambda\) of \(\operatorname{Gal}(\mathbf{N}/\mathbf{K})\) into \(\varprojlim \operatorname{Gal}(\mathbf{N}_i/\mathbf{K})\) is an isomorphism of topological groups.

Since \(\operatorname{Gal}(\mathbf{N}/\mathbf{K})\) is compact, \(\lambda\) continuous and the group \(\varprojlim \operatorname{Gal}(\mathbf{N}_i/\mathbf{K})\) separated, it suffices to show that \(\lambda\) is bijective (Gen.~Top., I, p.~87, Cor. 2). Let \(u = (u_i)_{i \in I}\) be an element of \(\varprojlim \operatorname{Gal}(\mathbf{N}_i/\mathbf{K})\) ; for each \(i \in I\) , \(u_i\) is a K-automorphism of \(\mathbf{N}_i\) and \(u_i\) is the restriction of \(u_j\) to \(\mathbf{N}_i\) for \(i \leqslant j\) . Since \(\mathbf{N} = \bigcup_{i \in I} \mathbf{N}_i\) there exists a unique element \(\sigma\) of \(\operatorname{Gal}(\mathbf{N}/\mathbf{K})\) which coincides with

\(u_{i}\) on \(\mathbf{N}_i\) for all \(i\in \mathbf{I}\) . Thus \(\sigma\) is the unique element of \(\operatorname {Gal}(\mathbf{N} / \mathbf{K})\) such that \(\lambda (\sigma) = u\) , hence \(\lambda\) is bijective.

This applies in particular when we take for the family \((\mathbf{N}_{i})\) the family of all finite Galois subextensions of N; then each group \(\operatorname{Gal}(\mathbf{N}_{i}/\mathbf{K})\) is discrete and finite. The topological group \(\operatorname{Gal}(\mathbf{N}/\mathbf{K})\) is thus isomorphic to a directed inverse limit of finite groups, equipped with the discrete topology; this is sometimes called a profinite topological group.

